\chapter{State Counting Mass Spectrometry}
\label{ch:state_counting}

\papernote{Source paper: \emph{State Counting Mass Spectrometry} --- \texttt{mekaneck/docs/state-counting/}}

\begin{chapterabstract}
State counting inverts the mass spectrometry paradigm from continuous analogue to fundamentally digital. Ion trajectories traverse discrete partition states $(n, \ell, m, s)$; measurement counts state traversals until trajectory completion at the $\varepsilon$-boundary. The state-mass correspondence theorem establishes a bijection between partition state count and $m/z$. Time-state identity: $\mathrm{d}M/\mathrm{d}t = 1/\langle\tau_p\rangle$. The result is a mass spectrometry instrument whose operation is a direct physical implementation of the categorical counting operation, providing a uniquely transparent experimental probe of the framework's foundational claims.
\end{chapterabstract}

\input{../mekaneck/docs/state-counting/state-counting-mass-spectrometry.tex}
